% Naslovna strana, KDI (sr/en), Zadatak, Izjava, epigraf.
% Rucno pisano (ne kroz konverter) - ove strane prate zvanican FTN obrazac i
% moraju sadrzati stvarne potpise/datume pri predaji, ne samo formatiran tekst.

\begin{titlepage}
\centering
{\bfseries УНИВЕРЗИТЕТ У НОВОМ САДУ\\FAKULTET TEHNIČKIH NAUKA U NOVOM SADU}\\[2pt]
\rule{\textwidth}{0.8pt}

\vspace{5cm}
\begin{flushleft}
\Autor
\end{flushleft}

\vspace{4.5cm}
{\LARGE\bfseries \Naslov}

\vspace{1cm}
Diplomski rad\\
- Osnovne akademske studije -

\vfill
Novi Sad, 2026.
\end{titlepage}

\chapter*{Ključna dokumentacijska informacija}
\begin{longtable}{@{}p{0.32\textwidth}p{0.6\textwidth}@{}}
\textbf{Redni broj, RBR:} & \\
\textbf{Identifikacioni broj, IBR:} & \\
\textbf{Tip dokumentacije, TD:} & Monografska dokumentacija \\
\textbf{Tip zapisa, TZ:} & Tekstualni štampani materijal \\
\textbf{Vrsta rada, VR:} & Diplomski rad \\
\textbf{Autor, AU:} & \Autor \\
\textbf{Mentor, MN:} & \Mentor \\
\textbf{Naslov rada, NR:} & \Naslov \\
\textbf{Jezik publikacije, JP:} & Srpski / latinica \\
\textbf{Jezik izvoda, JI:} & Srpski \\
\textbf{Zemlja publikovanja, ZP:} & Republika Srbija \\
\textbf{Uže geografsko područje, UGP:} & Vojvodina \\
\textbf{Godina, GO:} & 2026 \\
\textbf{Izdavač, IZ:} & Autorski reprint \\
\textbf{Mesto i adresa, MA:} & Novi Sad, Trg Dositeja Obradovića 6 \\
\textbf{Fizički opis rada, FO:} & 10/[BROJ]/17/[BROJ]/[BROJ]/3/0 \\
\textbf{Naučna oblast, NO:} & Elektrotehnika i računarstvo \\
\textbf{Naučna disciplina, ND:} & Primenjeno softversko inženjerstvo \\
\textbf{Predmetna odrednica/Ključne reči, PO:} & Rutiranje teretnih vozila, OpenStreetMap, Valhalla, distribuirani sistemi, mobilna aplikacija, prilagođena funkcija cene rute \\
\textbf{UDK} & \\
\textbf{Čuva se, ČU:} & U biblioteci Fakulteta tehničkih nauka, Novi Sad \\
\textbf{Važna napomena, VN:} & \\
\textbf{Izvod, IZ:} & U ovom radu opisan je sistem za rutiranje teretnih vozila koji uzima u obzir fizička ograničenja vozila (visina, širina, dužina, masa, osovinsko opterećenje, prevoz opasnog tereta) i podatke o putnoj mreži preuzete iz OpenStreetMap projekta. Sistem se sastoji od Go backend servisa, Valhalla routing engine-a sa truck costing modelom, PostgreSQL/PostGIS baze podataka, RabbitMQ mehanizma za asinhronu obradu i Flutter mobilne aplikacije za vozače i dispečere. Nad Valhalla alternativama implementiran je sopstveni modul za rangiranje ruta prema riziku i preferencama vozača, kao i mehanizam koji vozaču objašnjava zašto predložena ruta odstupa od geometrijski najkraćeg puta. Za potrebe algoritamske evaluacije implementirana je sopstvena Dijkstra/A* pretraga nad ograničenim podgrafom putne mreže. \\
\textbf{Datum prihvatanja teme, DP:} & \\
\textbf{Datum odbrane, DO:} & \\
\textbf{Članovi komisije, KO:} & Predsednik: \\
 & Član: \\
 & Član, mentor: \Mentor \\
\end{longtable}

\chapter*{Key words documentation}
\begin{longtable}{@{}p{0.32\textwidth}p{0.6\textwidth}@{}}
\textbf{Document type, DT:} & Monographic publication \\
\textbf{Contents code, CC:} & Bachelor thesis \\
\textbf{Author, AU:} & [NAME SURNAME] \\
\textbf{Mentor, MN:} & [MENTOR TITLE AND NAME] \\
\textbf{Title, TI:} & Development of a heavy vehicle routing application adapted to the physical constraints of the vehicle and terrain \\
\textbf{Language of text, LT:} & Serbian \\
\textbf{Country of publication, CP:} & Republic of Serbia \\
\textbf{Publication year, PY:} & 2026 \\
\textbf{Scientific field, SF:} & Electrical and Computer Engineering \\
\textbf{Subject/Key words, S/KW:} & Heavy vehicle routing, OpenStreetMap, Valhalla, distributed systems, mobile application, custom route cost function \\
\textbf{Abstract, AB:} & This thesis describes a heavy vehicle routing system that accounts for the vehicle's physical constraints using road network data from OpenStreetMap, combining the Valhalla routing engine's truck costing model with a custom risk-ranking and route-explanation layer, a Go backend, RabbitMQ-based asynchronous processing, and a Flutter mobile application for drivers and dispatchers. \\
\end{longtable}

\chapter*{Zadatak za izradu diplomskog (Bachelor) rada}
\textbf{Vrsta studija:} Osnovne akademske studije\\
\textbf{Studijski program:} [STUDIJSKI PROGRAM]\\
\textbf{Student:} \Autor \quad \textbf{Broj indeksa:} [BROJ INDEKSA]\\
\textbf{Oblast:} Elektrotehničko i računarsko inženjerstvo\\
\textbf{Mentor:} \Mentor

\bigskip
\textbf{Tekst zadatka:}
\begin{itemize}
\item Proučiti postojeća komercijalna i otvorena rešenja za rutiranje teretnih vozila i principe rada distribuiranih sistema.
\item Proučiti podatkovni model OpenStreetMap projekta u delu koji se odnosi na ograničenja za teretni saobraćaj i pripremiti geografske podatke za teritoriju Republike Srbije.
\item Implementirati backend servis koji generiše, rangira sopstvenom funkcijom rizika i objašnjava predloženu rutu za dato vozilo.
\item Implementirati, u svrhu algoritamske evaluacije, sopstvenu pretragu najkraćeg puta (Dijkstra/A*) nad ograničenim podgrafom putne mreže.
\item Implementirati distribuiranu komunikaciju sistema i mobilnu aplikaciju za vozače i dispečere.
\item Testirati i evaluirati implementirano rešenje i izvesti najvažnije zaključke.
\end{itemize}

\chapter*{Izjava o nepostojanju sukoba interesa}
Izjavljujem da nisam u sukobu interesa u odnosu mentor -- kandidat i da nisam član porodice, povezano lice, odnosno da nisam zavisan/na od mentora/kandidata, da ne postoje okolnosti koje bi mogle da utiču na moju nepristrasnost, niti da stičem bilo kakve koristi ili pogodnosti bilo pozitivnim ili negativnim ishodom.

\bigskip
U Novom Sadu, dana \underline{\hspace{4cm}}

\bigskip
Mentor \dotfill\\[6pt]
Kandidat \dotfill

\clearpage
\thispagestyle{empty}
\vspace*{0.35\textheight}
\begin{center}
\itshape „The map is not the territory.“\\[4pt]
\upshape --- Alfred Korzybski
\end{center}
\clearpage
